\section*{अध्याय \ref{ch:b1:series} --- संख्यात्मक श्रेणी}

\begin{solution}{exo:b1:series:1}
$\dfrac{1}{n(n+1)} = \dfrac1n - \dfrac{1}{n+1}$: दूरबीनी, $S_N = 1 -
\frac{1}{N+1} \to 1$। अभिसारी, योग $1$।

$\ln\bigl(1 - \frac{1}{n^2}\bigr) = \ln\frac{(n-1)(n+1)}{n^2} =
\ln\frac{n-1}{n} - \ln\frac{n}{n+1}$: फिर से दूरबीनी, $S_N = \ln\frac12 -
\ln\frac{N}{N+1} \to -\ln 2$। अभिसारी, योग $-\ln 2$।

$\dfrac{3^n + 4^n}{5^n} = \bigl(\frac35\bigr)^n + \bigl(\frac45\bigr)^n$: दो
अभिसारी गुणोत्तर श्रेणियाँ, योग $\frac{1}{1 - 3/5} + \frac{1}{1 - 4/5} =
\frac52 + 5 = \frac{15}{2}$।
\end{solution}

\begin{solution}{exo:b1:series:2}
सर्वत्र \omterm{thm:b1:series:ratio}{अनुपात कसौटी}।

$\frac{u_{n+1}}{u_n} = \frac{(n+1)^2}{2n^2} \to \frac12 < 1$: अभिसारी।

$\frac{u_{n+1}}{u_n} = \frac{(n+1)!\,n^n}{n!\,(n+1)^{n+1}} =
\bigl(\frac{n}{n+1}\bigr)^n = \bigl(1 + \frac1n\bigr)^{-n} \to \frac1\eu <
1$: अभिसारी।

गुणक $2^n$ के साथ: अनुपात $\to \frac2\eu < 1$: अभिसारी।

$3^n$ के साथ: अनुपात $\to \frac3\eu > 1$: अपसारी (पद $+\infty$ की ओर जाते
हैं)।
\end{solution}

\begin{solution}{exo:b1:series:3}
सभी पद ऋणेतर; तुल्यताओं का प्रयोग कीजिए (\cref{thm:b1:series:positive})।

$\sin\frac{1}{n^2} \sim \frac{1}{n^2}$: अभिसारी (रीमान $\alpha = 2$)।

$1 - \cos\frac1n \sim \frac{1}{2n^2}$: अभिसारी।

$\frac{1}{\sqrt{n(n+1)}} \sim \frac1n$: अपसारी।

$\frac{\ln n}{n^2} = \frac{1}{n^{3/2}}\cdot\frac{\ln n}{n^{1/2}}$
और $\frac{\ln n}{\sqrt n} \to 0$ (\cref{prop:b1:functions:powerrules}): अतः
बड़े $n$ के लिए $\frac{\ln n}{n^2} \leq \frac{1}{n^{3/2}}$: अभिसारी।
\end{solution}

\begin{solution}{exo:b1:series:4}
$n \geq 2$ के लिए: $\frac{1}{n^2} \leq \frac{1}{n(n-1)} = \frac{1}{n-1} -
\frac1n$। अतः
\[
\sum_{n=1}^{N} \frac{1}{n^2} \leq 1 + \sum_{n=2}^{N}
\Bigl(\frac{1}{n-1} - \frac1n\Bigr) = 1 + 1 - \frac1N < 2 :
\]
आंशिक योग वर्धमान और $2$ से परिबद्ध: अभिसरण (\cref{thm:b1:series:positive}),
योग $\leq 2$। (यथार्थ मान $\frac{\pi^2}{6}$ द्वितीय वर्ष का उत्सव है।)
\end{solution}

\begin{solution}{exo:b1:series:5}
$\sum \frac{(-1)^n}{\sqrt n}$: $\frac{1}{\sqrt n} \downarrow 0$ के साथ
एकांतरित: अभिसारी (\cref{thm:b1:series:alternating}); निरपेक्ष रूप से नहीं
($\alpha = \frac12 \leq 1$)।

$\sum \frac{(-1)^n}{n + (-1)^n}$: अनुक्रम $\frac{1}{n + (-1)^n}$ ह्रासमान
\emph{नहीं} है ($\frac{1}{n+1}$ फिर $\frac{1}{n}$ बुरी तरह एकांतरित होते
हैं), अतः कसौटी सीधे लागू नहीं होती। प्रसार कीजिए:
\[
\frac{(-1)^n}{n + (-1)^n}
= \frac{(-1)^n}{n}\cdot\frac{1}{1 + \frac{(-1)^n}{n}}
= \frac{(-1)^n}{n} - \frac{1}{n^2} +
O\Bigl(\frac{1}{n^3}\Bigr):
\]
पहली \omterm{def:b1:series:def}{श्रेणी} अभिसरित होती है (एकांतरित), $\sum \frac{1}{n^2}$ अभिसरित होती
है, $O(n^{-3})$ निरपेक्ष रूप से अभिसरित होती है: तीन अभिसारी श्रेणियों का
योग अभिसरित होता है।

$\sin\bigl(\pi\sqrt{n^2+1}\bigr)$: $\varepsilon_n = O(n^{-3})$ के साथ
$\sqrt{n^2 + 1} = n + \frac{1}{2n} + \varepsilon_n$ लिखिए; फिर, चिह्न तक
$\sin$ की $\pi$-आवर्तिता से,
\[
\sin\bigl(\pi\sqrt{n^2+1}\bigr)
= (-1)^n \sin\Bigl(\frac{\pi}{2n} + \pi\varepsilon_n\Bigr) .
\]
$\theta_n = \frac{\pi}{2n} + \pi\varepsilon_n$ और $a_n = \sin\theta_n$ रखिए।
बड़े $n$ के लिए $\theta_n \in \intoo{0}{\frac\pi2}$ और
\[
\theta_n - \theta_{n+1} = \frac{\pi}{2n(n+1)} +
\pi(\varepsilon_n - \varepsilon_{n+1})
= \frac{\pi}{2n^2} + O\Bigl(\frac{1}{n^3}\Bigr) > 0
\]
अंततः, अतः $(\theta_n)$ घटकर $0$ पर आता है; चूँकि $\sin$
$\intcc{0}{\frac\pi2}$ पर वर्धमान है, $(a_n)$ भी घटकर $0$ पर आता है।
एकांतरित कसौटी लागू होती है: अभिसारी --- निरपेक्ष रूप से नहीं, क्योंकि $a_n
\sim \frac{\pi}{2n}$।
\end{solution}

\begin{solution}{exo:b1:series:6}
$f(t) = \frac{1}{t(\ln t)^\alpha}$ $\intco{2}{+\infty}$ पर धनात्मक, \omterm{def:b1:continuity:continuous}{संतत} और
ह्रासमान है। $u = \ln t$ प्रतिस्थापित करने पर:
\[
\int_2^x \frac{\dd t}{t(\ln t)^\alpha}
= \int_{\ln 2}^{\ln x} \frac{\dd u}{u^\alpha},
\]
$x \to \infty$ होने पर परिबद्ध तभी जब $\alpha > 1$
(\cref{thm:b1:series:riemann} की संगणना)। समाकल-तुलना से: अभिसरण तभी जब
$\alpha > 1$। (ये \emph{बेर्त्रां-प्रकार} की श्रेणियाँ दिखाती हैं कि अभिसरण
की सीमा-रेखा कितनी महीन है: $n\ln n$ अपसरित होती है, $n(\ln n)^{1.01}$
अभिसरित।)
\end{solution}

\begin{solution}{exo:b1:series:7}
\cref{exo:b1:derivative:3} के स्पर्शरेखा-परिबंधों को प्रसारों के द्वारा फिर
से लिखने पर: $x = \frac1n \in \intoc{0}{1}$ के लिए, $\ln(1+x)$ हेतु कोटि $1$
पर टेलर--लाग्रांज किसी $c \in \intoo{0}{x}$ के लिए $\ln(1 + x) = x -
\frac{x^2}{2(1 + c)^2}$ देता है, अतः
\[
0 \leq u_n = \frac1n - \ln\Bigl(1 + \frac1n\Bigr) \leq
\frac{1}{2n^2} .
\]
\omterm{thm:b1:series:riemann}{रीमान श्रेणी} से तुलना: $\sum u_n$ अभिसरित होती है। उसका आंशिक योग लघुगणकों
को दूरबीनी कर देता है:
\[
\sum_{n=1}^{N} u_n = H_N - \ln(N+1)
\]
(क्योंकि $\sum_{n\leq N} \ln\frac{n+1}{n} = \ln(N+1)$)। अतः $H_N - \ln(N+1)$
अभिसरित होती है; $\ln\frac{N+1}{N} \to 0$ जोड़ने पर अनुक्रम $H_N - \ln N$
अभिसरित होता है। उसकी सीमा $\gamma \approx 0.5772$ है।
\end{solution}

\begin{solution}{exo:b1:series:8}
$\dfrac{1}{n(n+2)} = \dfrac{1/2}{n} - \dfrac{1/2}{n+2}$: आंशिक योग $2$ के
\omterm{prop:b1:reals:intervals}{अंतराल} से दूरबीनी होता है,
\[
S_N = \frac12\Bigl(1 + \frac12 - \frac{1}{N+1} -
\frac{1}{N+2}\Bigr) \longrightarrow \frac34 .
\]

$\sum \frac{n}{2^n}$: $\abs x < 1$ के लिए, परिमित गुणोत्तर योग का अवकलन करके
सीमा तक जाने पर (यहाँ सभी श्रेणियाँ निरपेक्ष रूप से अभिसरित होती हैं, \omterm{thm:b1:series:ratio}{अनुपात
कसौटी}): $\sum_{n\geq0} x^n = \frac{1}{1-x}$ से, आंशिक योगों के साथ सीधी
संगणना करने पर मिलता है
\[
\sum_{n=1}^{N} n x^{n-1}
= \frac{1 - (N+1)x^N + N x^{N+1}}{(1 - x)^2}
\xrightarrow[N\to\infty]{} \frac{1}{(1-x)^2}
\quad (\abs x < 1),
\]
(परिसीमा-पद $N x^N \to 0$)। $x = \frac12$ पर: $\sum_{n\geq1}
n\bigl(\frac12\bigr)^{n-1} = 4$, अतः $\sum_{n\geq0} \frac{n}{2^n} = \frac12
\times 4 = 2$।
\end{solution}

\begin{solution}{exo:b1:series:9}
$\sum u_n$ के पदों को $2$ की घातों के बीच पैकेटों में समूहबद्ध कीजिए। ऊपरी
पैकेट: $2^k \leq n < 2^{k+1}$ के लिए $2^k$ पद हैं, प्रत्येक $\leq u_{2^k}$:
\[
\sum_{n=1}^{2^{K+1}-1} u_n
= \sum_{k=0}^{K} \sum_{n=2^k}^{2^{k+1}-1} u_n
\leq \sum_{k=0}^{K} 2^k u_{2^k} .
\]
निचले पैकेट: उसी पैकेट का प्रत्येक पद $\geq u_{2^{k+1}}$ है, अतः
$\sum_{n=2^k}^{2^{k+1}-1} u_n \geq 2^k u_{2^{k+1}} = \frac12 \cdot 2^{k+1}
u_{2^{k+1}}$, जिससे
\[
\sum_{n=1}^{2^{K+1}-1} u_n \geq \frac12 \sum_{k=1}^{K+1} 2^{k}
u_{2^{k}} .
\]
आंशिक योगों की दोनों तुलनाएँ दोनों दिशाओं में चलती हैं (ऋणेतर पद,
\cref{thm:b1:series:positive}): दोनों श्रेणियों की प्रकृति एक ही है।

रीमान: $u_n = n^{-\alpha}$ $2^k u_{2^k} = 2^{k(1-\alpha)}$ देता है, जो
\omterm{ex:b1:series:geometric}{गुणोत्तर श्रेणी} है, अभिसारी तभी जब $2^{1 - \alpha} < 1$, अर्थात् तभी जब
$\alpha > 1$। बेर्त्रां (\cref{exo:b1:series:6}): $u_n = \frac{1}{n(\ln
n)^\alpha}$ $2^k u_{2^k} = \frac{1}{(k\ln 2)^\alpha}$ देता है, जो $k$ में
\omterm{thm:b1:series:riemann}{रीमान श्रेणी} है: अभिसारी तभी जब $\alpha > 1$।
\end{solution}

\begin{solution}{exo:b1:series:10}
अनुपात $-t^2$ वाली परिमित गुणोत्तर सर्वसमिका:
\[
\frac{1}{1 + t^2} = \sum_{k=0}^{n-1} (-1)^k t^{2k}
+ \frac{(-1)^n t^{2n}}{1 + t^2} .
\]
$\intcc{0}{1}$ पर समाकलन कीजिए (बायाँ पक्ष समाकलित होकर $\arctan 1 =
\frac\pi4$ देता है, \cref{prop:b1:functions:arcderiv}):
\[
\frac{\pi}{4} = \sum_{k=0}^{n-1} \frac{(-1)^k}{2k+1}
+ (-1)^n \int_0^1 \frac{t^{2n}}{1+t^2}\,\dd t,
\qquad
0 \leq \int_0^1 \frac{t^{2n}}{1+t^2}\,\dd t \leq \int_0^1 t^{2n}\dd
t = \frac{1}{2n+1} .
\]
$n \to \infty$ लेने पर सूत्र सिद्ध हो जाता है, और \omterm{thm:b1:integration:def}{समाकल} पर प्रदर्शित परिबंध
ठीक शेषफल-परिबंध है: $N$ तक योग करने के बाद (अर्थात् $n = N + 1$ पद),
$\abs{R_N} \leq \frac{1}{2N + 3}$।
\end{solution}

\begin{solution}{exo:b1:series:11}
$n^{1/n} = \eu^{\frac{\ln n}{n}} \to \eu^0 = 1$
(\cref{prop:b1:functions:powerrules})। अतः
\[
\frac{1}{n^{1 + 1/n}} = \frac{1}{n}\,\eu^{-\frac{\ln n}{n}}
\sim \frac{1}{n} ,
\]
और तुल्यता कसौटी (\cref{thm:b1:series:positive}) अपसारी हरात्मक \omterm{def:b1:series:def}{श्रेणी} से
तुलना करती है: \emph{अपसारी}, यद्यपि हर घातांक $1 + \frac1n$ $1$ से बड़ा है।
रीमान कसौटी किसी \emph{नियत} घातांक $\alpha$ की बात करती है; $1$ की ओर सरकता
हुआ घातांक अपना पूरा अंतर गँवा सकता है, जैसा यहाँ हुआ।
\end{solution}

\begin{solution}{exo:b1:series:12}
मान लीजिए $\varepsilon > 0$। अभिसारी \omterm{def:b1:series:def}{श्रेणी} के लिए कोशी कसौटी से (आंशिक
योगों पर लगाया गया \cref{thm:b1:seq:complete}), ऐसा $N$ है कि $n \geq N$ के
लिए $\sum_{k=n+1}^{2n} u_k \leq \varepsilon$। एकदिष्टता से इन $n$ पदों में
से प्रत्येक $\geq u_{2n}$ है:
\[
n\,u_{2n} \leq \sum_{k=n+1}^{2n} u_k \leq \varepsilon
\quad\Longrightarrow\quad 2n\,u_{2n} \leq 2\varepsilon ,
\]
और विषम सूचकांकों के लिए $n \geq N$ हेतु $(2n+1)\,u_{2n+1} \leq
(2n+1)\,u_{2n} \leq 2\bigl(2n\,u_{2n}\bigr) \leq 4\varepsilon$: दोनों समताओं
में, $n u_n \to 0$।

विलोम असत्य: $u_n = \frac{1}{n\ln n}$ में $n u_n = \frac{1}{\ln n} \to 0$
है, फिर भी \omterm{def:b1:series:def}{श्रेणी} अपसरित होती है (\cref{exo:b1:series:6}, $\alpha = 1$)।
एकदिष्टता आवश्यक: $n$ के पूर्ण वर्ग होने पर $u_n = \frac1n$ और अन्यथा $u_n =
2^{-n}$ लीजिए: \omterm{def:b1:series:def}{श्रेणी} अभिसरित होती है (वर्ग वाले पद $\sum \frac{1}{k^2}$ की
तरह जुड़ते हैं, शेष गुणोत्तर रूप से), पर वर्गों के अनुदिश $n u_n = 1$।
\end{solution}

\begin{solution}{pb:b1:series:1}
\textbf{1.} $a_{n+1} - a_n = \frac{1}{n+1} - \ln\frac{n+1}{n} \leq 0$
क्योंकि $\ln(1 + \frac1n) \geq \frac{1/n}{1 + 1/n} = \frac{1}{n+1}$; और
$b_{n+1} - b_n = \frac{1}{n+1} - \ln\frac{n+2}{n+1} \geq 0$ क्योंकि $\ln(1 +
\frac{1}{n+1}) \leq \frac{1}{n+1}$। उनका अंतर $a_n - b_n = \ln(1 + \frac1n)
\to 0$: संलग्न (\cref{thm:b1:seq:adjacent}), उभयनिष्ठ सीमा $\lim a_n =
\gamma$ (\cref{exo:b1:series:7}) के साथ, और $b_n \leq \gamma \leq a_n$।

\textbf{2.} $b_{10} = 2.928968 - \ln 11 = 0.5311$ और $a_{10} = 2.928968 -
\ln 10 = 0.6264$: अतः $\gamma \in \intcc{0.5311}{0.6264}$। अंतर $\ln 1.1
\approx 0.095$ है और $\frac1n$ की तरह सिकुड़ता है: चार दशमलव ($\text{अंतर}
\leq 10^{-4}$) के लिए $n \approx 10^4$ चाहिए होगा --- घेराव सही हैं पर धीमे।

\textbf{3.} $a_n - a_{m+1} = \sum_{k=n}^{m} w_k$ को दूरबीनी करके $m \to
\infty$ लेने पर: $a_n - \gamma = \sum_{k\geq n} w_k$ (आंशिक योगों की सीमा)।
इसके अतिरिक्त
\[
w_k = \int_k^{k+1} \frac{\dd t}{t} - \frac{1}{k+1}
= \int_k^{k+1} \Bigl(\frac1t - \frac{1}{k+1}\Bigr)\dd t
= \int_k^{k+1} \frac{k + 1 - t}{t\,(k+1)}\,\dd t .
\]
$\intcc{k}{k+1}$ पर: $\frac{1}{(k+1)^2} \leq \frac{1}{t(k+1)} \leq
\frac{1}{k(k+1)}$, और $\int_k^{k+1}(k + 1 - t)\dd t = \frac12$: अतः
$\frac{1}{2(k+1)^2} \leq w_k \leq \frac{1}{2k(k+1)}$।

\textbf{4.} ऊपरी: $\sum_{k \geq n} \frac{1}{2k(k+1)} = \frac12\sum_{k\geq
n}\bigl(\frac1k - \frac{1}{k+1}\bigr) = \frac{1}{2n}$ (दूरबीनी)। निचला:
$\frac{1}{2(k+1)^2} \geq \frac{1}{2(k+1)(k+2)}$, जिसका योग दूरबीनी होकर
$\frac{1}{2(n+1)}$ देता है। प्रश्न 3 के साथ:
\[
\frac{1}{2(n+1)} \leq H_n - \ln n - \gamma \leq \frac{1}{2n} .
\]

\textbf{5.} $\frac{1}{2n}$ घटाइए: $\gamma_n - \gamma \in
\intcc{\frac{1}{2(n+1)} - \frac{1}{2n}}{0} = \intcc{-\frac{1}{2n(n+1)}}{0}$:
सुधारा गया आकलन $O\bigl(\frac{1}{n^2}\bigr)$ तक यथार्थ है, और सदा नीचे से।

\textbf{6.} $0 \leq \gamma - \gamma_{100} \leq \frac{1}{20200} <
5\cdot10^{-5}$ के साथ $\gamma_{100} = 5.1873775 - \ln 100 - 0.005 =
0.5772073$: अतः $0.577207 \leq \gamma \leq 0.577257$, अर्थात् $\gamma =
0.5772 \pm 5\cdot10^{-5}$ (सच्चा मान $0.5772156\dots$) --- सौ पदों से चार
प्रमाणित दशमलव, जबकि प्रश्न 2 के लिए दस हज़ार पद चाहिए थे।

\textbf{7.} पहला लाभांश:
\[
H_{2m} - H_m = \Bigl(\ln 2m + \gamma + \frac{1}{4m}\Bigr)
- \Bigl(\ln m + \gamma + \frac{1}{2m}\Bigr) +
O\Bigl(\frac{1}{m^2}\Bigr)
= \ln 2 - \frac{1}{4m} + O\Bigl(\frac{1}{m^2}\Bigr).
\]
दूसरा: $2m$ तक के सम व्युत्क्रमों का योग $\frac12 H_m$ है, अतः
$\sum_{j=1}^{m}\frac{1}{2j-1} = H_{2m} - \frac12 H_m = \frac12 \ln m + \ln 2
+ \frac\gamma2 + o(1)$, और $\sum_{j=1}^m \frac{1}{2j} = \frac12\ln m +
\frac\gamma2 + o(1)$।

\textbf{8.} सम पदों को दो बार अलग करने पर:
$\sum_{k=1}^{2m}\frac{(-1)^{k-1}}{k} = H_{2m} - 2\cdot\frac12 H_m = H_{2m} -
H_m$। प्रश्न 7 से यह $\ln 2 - \frac{1}{4m} + O(m^{-2})$ के बराबर है: योग
$\ln 2$ है (\cref{ex:b1:series:ln2} फिर से) \emph{और साथ में} उसकी गति भी।

\textbf{9.} $N = 2m$ के लिए: $S - S_N = \frac{1}{4m} + O(m^{-2}) =
\frac{1}{2N} + O(N^{-2})$। $N = 2m + 1$ के लिए: $S_{N} = S_{2m} +
\frac{1}{2m+1}$, अतः
\[
S - S_N = \Bigl(\frac{1}{4m} - \frac{1}{2m+1}\Bigr) +
O\Bigl(\frac{1}{m^2}\Bigr)
= -\frac{1}{4m} + O\Bigl(\frac{1}{m^2}\Bigr)
= -\frac{1}{2N} + O\Bigl(\frac{1}{N^2}\Bigr).
\]
दोनों स्थितियों में $S - S_N \sim \frac{(-1)^N}{2N}$: अधिकतम-स्थिति परिबंध
$a_{N+1}$ से आधा, और चिह्न ज्ञात तथा एकांतरित।

\textbf{10.} औसत लेने से दोलन करता अग्र पद मर जाता है:
\[
S - \tilde S_N = \frac{(S - S_N) + (S - S_{N+1})}{2}
= \frac{(-1)^N}{2}\Bigl(\frac{1}{2N} - \frac{1}{2(N+1)}\Bigr)
+ O\Bigl(\frac{1}{N^2}\Bigr) = O\Bigl(\frac{1}{N^2}\Bigr).
\]
संख्यात्मक रूप से: $S_{10} = 0.645635$, $S_{11} = 0.736544$, $\tilde S_{10}
= 0.691089$, और $\ln 2 = 0.693147$: त्रुटि $4.8\cdot10^{-2}$ से गिरकर
$2.1\cdot10^{-3}$ हो जाती है --- एक जोड़, बीस गुना बेहतर।

\textbf{11.} एकांतरित त्रुटि हर चरण पर चिह्न बदलती है, अतः क्रमागत आंशिक योग
सीमा को दोनों ओर से घेरते हैं और उनका माध्य प्रथम-कोटि पद काट देता है; घेराव
की त्रुटि $H_n - \ln n - \gamma \approx \frac{1}{2n}$ का चिह्न अचर है, अतः
$n$ के अनुदिश कोई भी औसत उसे नहीं काट सकता। दोनों \emph{घेरावों} का औसत लेना
सचमुच काम आता है: $\frac{a_n + b_n}{2} = H_n - \ln\sqrt{n(n+1)}$, और चूँकि
$\ln\sqrt{n(n+1)} = \ln\bigl(n + \frac12\bigr) + O(n^{-2})$,
\[
H_n - \ln\Bigl(n + \frac12\Bigr)
= \Bigl(H_n - \ln n - \frac{1}{2n}\Bigr) + \frac{1}{8n^2} +
O\Bigl(\frac{1}{n^3}\Bigr) = \gamma + O\Bigl(\frac{1}{n^2}\Bigr)
\]
(प्रश्न 5 और $\ln(1 + \frac{1}{2n}) = \frac{1}{2n} - \frac{1}{8n^2} +
O(n^{-3})$)। $n = 10$ पर जाँच: $H_{10} - \ln 10.5 = 0.57759$, जो पहले से ही
$\gamma$ के $4\cdot10^{-4}$ के भीतर है।

\textbf{12.} $\sum_{j\leq m} \frac{1}{2j-1} \geq \sum_{j \leq m}
\frac{1}{2j} = \frac12 H_m \to +\infty$: एकांतरित हरात्मक \omterm{def:b1:series:def}{श्रेणी} के धनात्मक
और ऋणात्मक दोनों भाग अपसरित होते हैं।

\textbf{13.} $u_n^\pm = \frac{\abs{u_n} \pm u_n}{2} \geq 0$ लिखिए, अतः $u_n
= u_n^+ - u_n^-$ और $\abs{u_n} = u_n^+ + u_n^-$। यदि $\sum u_n^+$ अभिसरित
होती, तो $\sum u_n^- = \sum (u_n^+ - u_n)$ भी अभिसरित होती (अभिसारी
श्रेणियों का अंतर), अतः $\sum \abs{u_n}$ भी: सप्रतिबंध अभिसरण से विरोधाभास।
सममिति से दोनों $\sum u_n^\pm$ अपसरित होती हैं ($+\infty$ की ओर): धनात्मक और
ऋणात्मक द्रव्यमान का अनंत भंडार।

\textbf{14.} मान लीजिए $S = \sum u_n$, $\varepsilon > 0$, और $\sum_{n > N}
\abs{u_n} \leq \varepsilon$ के साथ $N$ ($\sum\abs{u_n}$ के लिए कोशी कसौटी)।
$M_0$ इतना बड़ा लीजिए कि $\sigma(\{0, \dots, M_0\}) \supseteq \{0, \dots,
N\}$। $M \geq M_0$ के लिए, अंतर $\sum_{m \leq M} u_{\sigma(m)} - \sum_{n
\leq N} u_n$ \emph{भिन्न-भिन्न} पदों $u_n$ का परिमित योग है जिनमें $n > N$,
अतः उसका निरपेक्ष मान $\leq \varepsilon$; और $\abs{S - \sum_{n\leq N} u_n}
\leq \varepsilon$ भी। अतः पुनर्क्रमित आंशिक योग अंततः $S$ के $2\varepsilon$
के भीतर रहते हैं: $\sum u_{\sigma(n)} = S$। \omterm{thm:b1:series:absolute}{निरपेक्ष अभिसरण} पुनर्क्रमण-रोधी
है।

\textbf{15.} लालची प्रक्रिया का हर चरण परिमित कितने पदों के बाद समाप्त हो
जाता है, क्योंकि शेष धनात्मक (क्रमशः ऋणात्मक) पद अकेले ही अपसारी आंशिक योग
रखते हैं (प्रश्न 12): चलता योग अंततः $t$ पार कर ही लेगा। अतः प्रक्रिया अनंत
कितने परिमित चरणों के बीच एकांतरित होती है, धनात्मक पदों को क्रम से और
ऋणात्मक पदों को क्रम से खाती हुई: हर पद ठीक एक बार प्रयुक्त होता है ---
अर्थात् एक पुनर्क्रमण। पहली बार पार करने के बाद, दो क्रमागत पार करने के बीच
आंशिक योग एकदिष्ट रूप से $t$ की ओर बढ़ते हैं, और पार करते समय वे अधिक से
अधिक अभी-अभी जोड़े गए पद जितना आगे निकल जाते हैं; चूँकि $j$-वें पार पर
प्रयुक्त पदों का सूचकांक अपनी \omterm{def:b1:series:def}{श्रेणी} में कम से कम $j$ है, ये अतिक्रमण $0$ की
ओर जाते हैं। अतः आंशिक योग $t$ की ओर अभिसरित होते हैं: हर वास्तविक संख्या
किसी न किसी पुनर्क्रमण का योग है।

\textbf{16.} धनात्मक स्थानों को क्रम से $j = 1, 2, \dots$ हेतु
$\frac{1}{2j-1}$ मिलते हैं, ऋणात्मक स्थानों को क्रम से $\frac{1}{2j}$:
एकांतरित हरात्मक \omterm{def:b1:series:def}{श्रेणी} का हर पद ठीक एक बार आता है। $(p, q) = (1, 2)$ के लिए
खंड $\bigl(1, -\frac12, -\frac14\bigr)$, $\bigl(\frac13, -\frac16,
-\frac18\bigr)$, $\bigl(\frac15, -\frac1{10}, -\frac1{12}\bigr)$, \dots\ हैं
--- अर्थात् प्रदर्शित \omterm{def:b1:series:def}{श्रेणी}।

\textbf{17.} चूँकि $\frac{1}{4k-2} = \frac12\cdot\frac{1}{2k-1}$:
\[
\frac{1}{2k-1} - \frac{1}{4k-2} - \frac{1}{4k}
= \frac12\,\frac{1}{2k-1} - \frac12\,\frac{1}{2k}
= \frac12\Bigl(\frac{1}{2k-1} - \frac{1}{2k}\Bigr).
\]
$k = 1, \dots, K$ पर योग करने पर: $T_{3K} = \frac12
\sum_{k=1}^{K}\bigl(\frac{1}{2k-1} - \frac{1}{2k}\bigr) = \frac12 S_{2K}$:
हर तीसरे आंशिक योग पर पुनर्क्रमित \omterm{def:b1:series:def}{श्रेणी} मूल \omterm{def:b1:series:def}{श्रेणी} की \emph{ठीक} आधी है।

\textbf{18.} $K$ पूरे खंडों के बाद पुनर्क्रमित आंशिक योग
$\sum_{j=1}^{pK}\frac{1}{2j-1} - \sum_{j=1}^{qK}\frac{1}{2j}$ है, और प्रश्न
7 उसका मूल्यांकन कर देता है:
\[
\Bigl(\frac{\ln(pK)}{2} + \ln 2 + \frac\gamma2\Bigr)
- \Bigl(\frac{\ln(qK)}{2} + \frac\gamma2\Bigr) + o(1)
= \ln 2 + \frac12\ln\frac pq + o(1) :
\]
$\gamma$ कट जाते हैं, $\ln K$ कट जाते हैं, और अनुपात $\frac pq$ बचा रहता है।

\textbf{19.} खंड $K + 1$ के भीतर का कोई आंशिक योग $K$-खंड वाले योग से अधिक
से अधिक $p + q$ पदों जितना भिन्न है, जिनमें से प्रत्येक का निरपेक्ष मान
$\leq \frac{1}{2qK}$ के आसपास है, अतः अंतर $O\bigl(\frac1K\bigr) \to 0$ है:
आंशिक योगों के पूरे अनुक्रम की सीमा वही $\ln 2 + \frac12\ln\frac pq$ है।
$(1, 2)$ के लिए: $\ln 2 + \frac12\ln\frac12 = \frac{\ln 2}{2} =
0.34657\dots$, और सचमुच $T_9 = 0.30833$ उसी की ओर रेंग रहा है: प्रश्न 17 से
$T_{3K} = \frac12 S_{2K}$ ठीक एकांतरित-हरात्मक त्रुटि की आधी त्रुटि के साथ
अभिसरित होता है। वही पद, आधा योग।

\textbf{20.} $(1,1)$: $\ln 2 + \frac12\ln 1 = \ln 2$ --- मूल क्रम, अर्थात्
संगति। $(2,1)$: $\frac32\ln 2 \approx 1.0397$। $(p,q)$ वाला मेनू ठीक गणनीय
\omterm{def:b1:topology:dense}{सघन} कुल $\ln 2 + \frac12\ln r$, $r \in \Q_{>0}$ तक पहुँचता है; रीमान की
लालची विधि (प्रश्न 15) \emph{हर} वास्तविक संख्या तक पहुँचती है। संरचना सूत्र
खरीदती है; लालच संपूर्णता।

\textbf{21.} आंशिक योग दूरबीनी हो जाते हैं: $\sum_{k=1}^{N} \bigl(\frac1k -
\ln\frac{k+1}{k}\bigr) = H_N - \ln(N+1) = b_N \to \gamma$:
\cref{exo:b1:series:7} की \omterm{def:b1:series:def}{श्रेणी} का योग ठीक \omterm{pb:b1:series:1}{ऑयलर का अचर} है।

\textbf{22.} $t = 1$ के लिए लालची विधि: पहला धनात्मक पद योग को ठीक $1$ तक
लाता है, उससे आगे नहीं, अतः पार करने के लिए दूसरा धनात्मक पद लिया जाता है:
$1, \frac13$ (योग $1.3333 > 1$), फिर $-\frac12$ ($0.8333$), $\frac15$
($1.0333$), $-\frac14$ ($0.7833$), $\frac17, \frac19$ ($1.0373$), $-\frac16$
($0.8706$), $\frac1{11}, \frac1{13}$ ($1.0384$), $-\frac18$ ($0.9134$),
$\frac1{15}$ ($0.9801$), \dots\ --- योग $1$ के इर्द-गिर्द निरंतर छोटे होते
आयाम के साथ साँस लेते हैं, और अब हर चक्र में दो धनात्मक पद चाहिए, क्योंकि
ऋणात्मक पद बड़े हैं।

\textbf{23.} खंड बनाइए: चरण $j$ पर इतने अप्रयुक्त धनात्मक पद जोड़िए कि आंशिक
योग कम से कम $1$ बढ़ जाए (यह संभव है: शेष धनात्मक पदों के योग अपसारी हैं,
प्रश्न 12), फिर अकेला ऋणात्मक पद $-\frac{1} {2j}$ जोड़िए। हर धनात्मक पद
अंततः प्रयुक्त होता है (हर चरण कम से कम एक लेता है), हर ऋणात्मक पद भी (प्रति
चरण एक): अर्थात् एक पुनर्क्रमण। हर चरण योग को $\geq 1 - \frac{1}{2j} \geq
\frac12$ जितना बदलता है: चरण $j$ के बाद आंशिक योग $\frac{j}{2}$ से अधिक हो
जाते हैं और किसी चरण के भीतर की वृद्धियाँ अंतिम को छोड़कर धनात्मक हैं, और
अंतिम $\frac{1}{2j} \to 0$ से परिबद्ध है: अतः $+\infty$ की ओर अपसरण।

\textbf{24.} नियम से $H_N \geq 20 \iff \ln N \geq 20 - \gamma - \frac{1}{2N}
+ O(N^{-2})$: देहली $N^*$ $\ln N^* = 20 - \gamma + o(1)$ संतुष्ट करती है,
अर्थात् $N^* = \eu^{20 - \gamma}(1 + o(1)) \approx \eu^{19.4228} \approx
2.72\cdot10^{8}$ --- जो \cref{ex:b1:series:harmonicstack} की स्थूल खिड़की
$\intcc{1.8\cdot10^8}{4.9\cdot10^8}$ के भीतर है, और $\gamma$ ने उसे कील दिया
है।

\textbf{25.} (क) $H_n = \ln n + \gamma + \frac{1}{2n} + O(n^{-2})$ में: $\ln
n$ \omterm{thm:b1:integration:def}{समाकल} है, $\gamma$ योग को \omterm{thm:b1:integration:def}{समाकल} से बदलने की क़ीमत (विश्लेषण का सचमुच नया
अचर), और $\frac{1}{2n}$ पहला सुधार --- समलंब की छाया। (ख) सप्रतिबंध अभिसरण
दो अनंत भंडारों के बीच के निरसन पर टिका है (प्रश्न 13), अतः क्रम बदलने से
भंडारों का भार बदल जाता है; \omterm{thm:b1:series:absolute}{निरपेक्ष अभिसरण} का कुल द्रव्यमान परिमित है, और
प्रश्न 14 का पूँछ-आकलन क्रम को देखता ही नहीं। (ग) $(p,q)$ वाले योग
\emph{संगणित} किए गए थे: $\gamma$-नियम ने हर पुनर्क्रमित आंशिक योग को $\ln 2
+ \frac12\ln\frac pq + o(1)$ में बदल दिया, जिसमें $\gamma$ स्वयं कट गया ---
यह अनंतस्पर्शी लेखा-जोखा था, कोई अमूर्त तर्क नहीं। (घ) आगे: स्नातक वर्ष 2 के
खंड में घात \omterm{def:b1:series:def}{श्रेणी} के लिए गुणनफल और अप्रतिबंध योगनीयता; और स्नातक वर्ष 3 के
खंड की स्टर्लिंग सूत्र वाली सप्ताहांत समस्या, जहाँ योग-बनाम-समाकल का वही
लेखा-जोखा एक कोटि आगे बढ़ाने पर स्वयं $\sqrt{2\pi}$ उत्पन्न कर देता है।
\end{solution}
